% SOURCE_AUTHORITY: sga5_fr_workpass.tex, lines 14040--14050 (LF-normalized unit).
% SOURCE_SHA256: 791F4EFFC5E02832D5D77ED03518C8156D6F07E4C8238B03545DB93D883FBB28
\subsubsection*{c) O caso «moderadamente ramificado».} Suponha-se que exista $x' \in p^{-1}(x)$ tal que
\[
f'(x') = x' \text{ , } \quad (p(x') = x) \quad ,
\]
caso ao qual sempre se pode reduzir (cf. o caso a)), e que a ordem do grupo $G_{x'}$ seja prima com a característica do corpo $k$. Neste caso, tem-se:

\begin{equation}\tag{6.4.4}
\alpha_x(f,u, f', G, \varphi) = S_{x,\ell^n}(f, f', G, \varphi)(e) \, \Tr_A(u_M^\vee) \quad .
\end{equation}

Levando em conta o caso b), pode-se supor que $G = G_{x'}$. A fórmula (6.4.4) resulta então do lema seguinte:
